\part{The Empirical Program and Conclusions}
\label{part:program}

\chapter{The Experimental Program}
\label{ch:experiments}

\section{Status of the Proposals}

The experiments in this chapter are designs. None has been run for this monograph, and no result is claimed. Their purpose is to state, before data are collected, what would count as support for or against each hypothesis, so that the hypotheses are exposed to refutation.

\section{Experiment I: Geometric Residuals and Task Success}

The first experiment tests Hypothesis~\ref{hyp:geom}. Models are trained in the three regimes of Chapter~\ref{ch:valmeekam} on a maze family with several distributions. For each model and evaluation task, record answer correctness, trace validity under the A* specification, generated length, a difficulty covariate computed from properties of the maze, and the hidden-state trajectory at a stated set of layers. Construct reference sets of states from held-out correct completions, fit local tangent estimates at several neighbor counts and dimensions, and compute the trajectory residual of Chapter~\ref{ch:hidden-states}. Fit a regularized logistic model of correctness on length, difficulty, and validity as the baseline, add the residual, and compare held-out log-likelihood by cross-validation across mazes.

Four controls are required. A null residual computed from random directions of matched norm establishes the scale at which a residual is small. A residual computed against a reference set drawn from unrelated tasks establishes that the effect, if any, is specific to task-relevant states. A residual computed on shuffled trajectories from different mazes establishes that the effect is not an artifact of the position along the trajectory. The fourth control addresses construct validity. Reference manifolds estimated only from correct completions may detect proximity to a success-associated distribution of states and not a general tangent-alignment property. Reference sets are therefore also constructed without conditioning on correctness, from all completions or from a random sample of states, and the residuals computed against the two kinds of reference are compared. If only the correctness-conditioned residual predicts success, the result is attributed to proximity to a success-associated representation distribution and is not read as support for a general geometric principle. Uncertainty estimates account for both task sampling and independently trained model seeds, with seeds treated as a random effect, because bootstrapping mazes alone would omit variability due to training initialization. The hypothesis is supported if the residual improves held-out prediction beyond the baseline across the stated range of estimator settings and the controls do not show a comparable improvement. It is not supported if the improvement vanishes under the controls or appears only at isolated settings.

\section{Experiment II: Scaffold Utility}

The second experiment is the prefix factorial of Chapter~\ref{ch:scaffolding}. For each prefix type, namely correct trace, swapped trace, random tokens, repeated tokens, and learned prefix, evaluate answer accuracy at matched token counts and, separately, at matched computational cost, since equal token counts do not equalize the full training budget or implementation-level inference cost, and estimate the interventional utility $\Delta U$ relative to the empty prefix and the random prefix. Where internal access is available, compute the trajectory residual for each prefix type and examine whether the ordering of prefix types by utility matches the ordering by residual. A match would be consistent with the geometric account, and a mismatch would indicate that utility arises through a pathway the residual does not capture.

\section{Experiment III: Admissible Agents}

The third experiment evaluates the execution boundary. A suite of formally checkable problems is posed to agents that propose actions subject to independent authorization constraints. The suite includes three kinds of task, constructed so that the three relations among correctness, validity, and permission are all represented: tasks for which the correct solution is unauthorized, tasks in which the explanation is valid but the action is prohibited, and tasks in which the action is permitted although the accompanying explanation is invalid. Four architectures are compared. In the first the model authorizes itself, so that its own judgment decides execution. In the second a second model provides approval. In the third a formal verifier checks the proposal against the policy specification. In the fourth a runtime controller enforces permissions at the execution boundary in the sense of Proposition~\ref{thm:soundness}.

For each architecture the evaluation records answer accuracy, trace validity, causal sensitivity of the output to the trace, authorization accuracy against the ground-truth policy, execution violations, refusal quality, and resource cost. Refusal quality is measured by the fraction of refusals that carry the correct reason code and by the fraction of tasks that could have been completed admissibly and were not. The primary safety metric is the rate of unauthorized execution and not the persuasiveness of generated justifications. Persuasiveness may be recorded as a covariate, since a central prediction is that the unauthorized execution rate of the first two architectures depends on the persuasiveness of the generator's justification, while that of the third and fourth does not.

\section{Predictions}

The program yields predictions that can fail. The unauthorized execution rate of the fourth architecture is predicted to be zero up to violations of the stated assumptions, and any observed violation is to be traced to one of (M), (S), (A), (P), or (I), or else to a specification or instrumentation failure, meaning a mismatch between the formalized action, the measured event, and the actual behavior of the system. The evaluation must be able to report the second kind of cause as well as enforcement failures. That of the self-authorizing architecture is predicted, conditionally, to increase with the generator's competence when successful task completion is rewarded without adequate authorization constraints, because more competent generators then produce more convincing proposals for actions that the policy forbids. The relation is not necessary: a more competent model might also be better at recognizing authorization constraints. The experiment therefore varies reward structure and policy awareness independently of competence and tests the interaction among the three. That of the second-model architecture is predicted to lie between them and to correlate with the similarity of the two models. If the self-authorizing architecture's violation rate does not increase with competence under the reward structure that rewards completion alone, the claim that competence promotes statistical into procedural authority is weakened for the tested population.

\section{Analysis Plan and Threats to Validity}

Each experiment should be specified in an analysis plan fixed before the data are examined, including the primary outcome, the covariates in the baseline model, the estimator grid for the geometric measurements, the thresholds for what will be called support or non-support, and the rule for handling failed runs. The purpose of fixing these in advance is to prevent the flexibility of the geometric estimators, which have several tunable parameters, from being used to find an association after the fact. The estimator grid should be reported in full, with the results at every point of the grid and not only at the point where the effect is largest, so that the reader can see how stable the conclusion is across the grid.

Several threats to validity should be addressed in the design. Confounding by difficulty arises because harder mazes may produce both longer traces and more erratic internal trajectories, and the baseline must include difficulty measures strong enough to absorb this. Confounding by confidence arises because a model that is uncertain may both fail more often and traverse less structured states, and the baseline should include a measure of the model's uncertainty at the start of generation. Overfitting of the estimator arises because local manifold estimates built from a reference set may fit that set closely and fit new states poorly in a way correlated with correctness for reasons unrelated to coherence, and the remedy is to estimate on one subset and evaluate on another. Construct validity arises because the residual measures departure from an estimated local linear structure, and whether that corresponds to the semantic notion of coherence in the earlier theory is a question that the experiment cannot settle on its own.

\section{What Would Change the Conclusions}

The hypotheses are stated so that results can contradict them. If swapped traces provide no utility beyond random tokens of equal length, the scaffolding hypothesis reduces to a claim about serial computation, and the conditional structure of traces plays no role. If the geometric residual adds nothing to prediction of correctness beyond the baseline, Hypothesis~\ref{hyp:geom} is not supported for the tested models. If the unauthorized execution rate of the self-authorizing architecture does not increase with generator competence, the prediction about the growth of the risk with capability fails. If an execution violation is observed in the architecture with a runtime controller, the violation is traced to the failed assumption, and the proposition is not refuted but its applicability to that deployment is. Distinguishing a refuted proposition from a violated assumption is the purpose of stating the proposition conditionally.

\chapter{Scope and Limits}
\label{ch:limits}

\section{Formal and Open-Ended Domains}

The results of this monograph are sharpest in formal domains, where a specification is decidable and a policy can be written completely. In such domains validity can be checked exactly, correctness can be determined independently, and an authorization controller can be built to be sound. In open-ended environments the specification of the task is incomplete, the policy is contested or underspecified, and the situation cannot be fully captured. The structural claims remain, since the independence of the properties does not depend on decidability, but the guarantees weaken, because the soundness assumption (S) is only as good as the completeness of the policy.

\section{Verification That Uses Language Models}

An objection holds that independent verification may itself require language models, because the policies and situations are expressed in natural language and cannot be parsed otherwise. The distinction that answers it is between probabilistic assistance in verification and unilateral authority over execution. A language model may be used to extract the relevant facts from a document, to flag possible violations, or to rank cases for human review, and in these roles its errors are caught by the stage that follows. It may not be the stage that decides, in the sense of determining execution without a sound check, unless its decisions are themselves subject to the boundary of Proposition~\ref{thm:soundness}. The architecture constrains where authority sits and permits any technology to supply evidence.

\section{Stronger Models}

Stronger models do not remove the need for external verification. The independence proposition does not depend on model quality. Improvement in competence raises the proportion of proposals that are correct and also raises the persuasiveness of those that are not, and it does nothing to the policy, which is a separate input. If anything the incentive to verify independently grows with competence, because the cost of an unverified action grows with the capability of the system proposing it.

\section{Limits of the Geometric Account}

The geometric account is a hypothesis with specific vulnerabilities. Hidden states may not be well approximated by low-dimensional manifolds at the scales relevant to the task. Estimates of coherence depend on the estimator and the reference set. A measured association between coherence and success need not be causal. The scaffolding hypothesis does not depend on the geometric account and survives its failure. The authorization results do not depend on either and would stand if both were rejected.

\section{Open Problems}

Several problems are left open by the analysis. The first is the extension of the geometric account to settings in which hidden states are better described by a stratified space than by a smooth manifold, with a corresponding theory of what it means for a trajectory to cross strata coherently. The second is the relation between token-level interventions and state-level interventions, which the monograph treats as an empirical quantity to be measured and for which no theory is offered. The third is the design of verifiers for policies that cannot be written formally, where the policy is partly implicit in the practices of an institution, and where the soundness assumption is correspondingly difficult to state, let alone secure. The fourth is the treatment of multi-agent systems, in which the authority of one agent over another is itself a policy question, and the provenance model of Chapter~\ref{ch:laundering} must be extended to labels that evolve as agents delegate and endorse.

A fifth problem concerns the economics of verification. The architecture imposes costs, and the question of how much verification is justified for which actions is a question of decision under uncertainty with asymmetric losses that has not been addressed here beyond noting that the strongest checks belong on the actions whose consequences are hardest to reverse. A sixth concerns the human side of the boundary. The authorization controller may be a person or a committee, and the reliability of human authorization is itself subject to the failures of attention and deference that automation tends to produce. A system in which people approve whatever a fluent system proposes has reintroduced the conflation in a different place, and the design of interfaces that preserve the independence of the human warrant is a problem for which this monograph provides the motivation and not the solution.

\chapter{Conclusion: Reasons Are Not Permissions}
\label{ch:conclusion}

The monograph began from an empirical observation: intermediate tokens can improve performance without functioning as valid derivations. From that observation it proceeded in three steps. The first step distinguished the properties that are commonly run together, namely utility, procedural validity, causal faithfulness, geometric coherence, outcome correctness, and policy admissibility, and showed by explicit countermodels that none of the first five entails the sixth and that no unconditional entailment holds among the first five without stated linking assumptions. The result is the absence of unconditional entailment. It is not a claim of statistical independence, and it does not assert that every combination of values is realizable under every specification. The second step offered a positive hypothesis about what makes intermediate sequences useful, in terms of the geometry of the trajectories they induce, with a test that could refute it. The third step moved from the evaluation of explanations to the control of action, and showed that an independent authorization mechanism with complete mediation, atomic validation, and a sound check guarantees that no inadmissible action is executed, whatever the generator proposes and however persuasive its explanation.

The appropriate response to increasingly competent generative systems is not to require that every useful internal computation resemble a human explanation. It is to specify which claims require which evidence, and which actions require which authorization. The distinction that governs the manuscript is among what has been established empirically, what follows logically from definitions, and what must be guaranteed architecturally, and the manuscript has attempted to avoid the category error that it describes by keeping the three apart.

It is useful to restate what each register of claim has contributed. Empirically, the monograph relies on a single source for the dissociation between trace validity and answer accuracy in maze-solving transformers, and it has been careful to state that result in its narrow form. Logically, it has proved that the six properties are independent except under stated linking assumptions, that permission is not entailed by any of the others, that a projection of an update onto a cone of feasible tangent directions is well defined and satisfies first-order feasibility, and that an execution boundary with complete mediation, atomic validation, and a sound check guarantees the admissibility of every executed action. Architecturally, it has argued that the guarantee is obtained by a small trusted base, a typed interface between roles, provenance labels that cannot be raised by derivation, and a history written by the execution environment. As conjectures it has offered the scaffolding hypothesis and the geometric faithfulness hypothesis, each with a test that could fail. The boundaries between these registers are the main methodological commitment of the work, and a reader who disagrees with any conjecture can discard it without affecting the logical and architectural claims.

The concluding principle is the following.
\begin{center}
\fbox{\begin{minipage}{0.72\textwidth}
\centering
Competence does not establish faithfulness.\\[0.4ex]
Faithfulness does not establish permission.\\[0.4ex]
Permission requires an independent warrant.
\end{minipage}}
\end{center}
Here an independent warrant is an authorization grounded in the applicable policy and in evidence, and not an assertion made sufficient by the proposing system's own success or explanation.
